\section{Objetivo general}

El objetivo de este trabajo es diseñar, describir y verificar un
sistema digital capaz de recibir dos operandos y un código de
operación mediante UART, ejecutar sobre ellos la ALU desarrollada en
el TP1 y transmitir el resultado por el mismo enlace. Además, el
sistema debe ser sintetizable, reproducible y capaz de volver al
estado inicial del protocolo después de recibir una trama incompleta o inválida.

\section{Objetivos específicos}

\begin{itemize}
	\item Mantener el contrato funcional de la ALU del TP1, incluidas las
		operaciones ADD, SUB, AND, OR, XOR, NOR, SRA y SRL.
	\item Implementar una UART 8N1 a \num{19200} baud con sobremuestreo 16x.
	\item Describir FSM separadas para re\-cep\-ción, transmisión y control
		de la tran\-sac\-ción con recuperación segura.
	\item Desacoplar comunicación y procesamiento mediante FIFOs.
	\item Definir un protocolo binario sobre el orden de los bytes y el
		tratamiento de errores.
	\item Verificar cada bloque de manera aislada y comprobar después el recorrido
		completo desde el pin RX hasta el pin TX.
	\item Proveer un cliente de Python que permita ejecutar la validación física.
\end{itemize}
